\chapter{Conclusioni}
\label{ch:conclusioni}

Le sette domande di ricerca formulate nel \Cref{ch:impianto} hanno trovato
risposta, una per una, nel \Cref{ch:valutazione}. Questo capitolo compone
quelle risposte in un bilancio conclusivo: sintetizza ciò che la tesi
dimostra e i contributi che ne derivano, discute le limitazioni del
prodotto, distinte da quelle della misura, e indica le direzioni lungo le
quali il lavoro può proseguire.

\section{Sintesi dei risultati e contributi}
\label{sec:concl-sintesi}

Il lavoro descritto in questa tesi ha mostrato che un sistema di supporto alla decisione clinica per la verifica della rimborsabilità delle Note AIFA (Agenzia Italiana del Farmaco) può essere costruito secondo una strategia neuro-simbolica praticabile, in cui la decisione di rimborsabilità è interamente affidata a un rule engine deterministico ispezionabile e la sola spiegazione clinica è prodotta da un large language model di taglia accessibile inserito in una pipeline retrieval-augmented. La separazione fra i due layer non è un dettaglio implementativo ma il cuore concettuale dell'approccio: rende il sistema verificabile per costruzione sul piano decisionale e, contemporaneamente, recupera dal componente generativo la capacità di articolare in linguaggio naturale ciò che il rule engine, da solo, riuscirebbe soltanto a sentenziare.

I risultati riportati nel \Cref{ch:valutazione} forniscono una risposta complessivamente positiva alle sette domande di ricerca. Il rule engine raggiunge una Macro F1 pari a $1{,}0000$ sui \goldn\ casi di riferimento, un valore atteso per costruzione che certifica la coerenza interna fra codifica delle regole e annotazione, entro i limiti di auto-referenza discussi in \Cref{sec:limitazioni}; il retrieval ha un Recall@10 di $0{,}9068$ e un MRR (\emph{Mean Reciprocal Rank}) di $1{,}0000$; la pipeline LLM produce spiegazioni con citation coverage del $100\%$ e hallucination rate nullo; il composito ALES (\emph{Aggregated LLM Evaluation Score}) vale $0{,}6810$. Sulla decisione di rimborsabilità il sistema completo distanzia di oltre il $68\%$ la baseline LLM con retrieval, e questo gap non è un miglioramento parametrico ma un effetto strutturale del rule engine deterministico. I test di idempotenza e di robustezza ai casi-frontiera restituiscono pass rate del $100\%$, in linea con la natura deterministica del layer simbolico. Le due risposte parziali (RQ3 sulla fedeltà letterale dell'LLM e RQ6 sulla verificabilità misurata da QAEval) non indicano debolezze del sistema ma limiti dei framework di valutazione automatica della fedeltà semantica e della copertura per micro-domande, come discusso in \Cref{sec:limitazioni}.

Il contributo della tesi si articola lungo tre piani distinti. Sul piano architetturale, la separazione fra layer simbolico ed esplicativo è una scelta di design replicabile su altre famiglie di Note AIFA e, più in generale, su altri sistemi prescrittivi regolati da specifiche testuali; non richiede modelli LLM proprietari, dataset di training annotati o infrastruttura specializzata oltre a un singolo nodo GPU per l'LLM locale. Sul piano metodologico, la valutazione strutturata per domande di ricerca supera l'antipattern dell'accumulazione indiscriminata di metriche e produce un messaggio sperimentale più leggibile per il revisore clinico e il valutatore accademico. Sul piano operativo, l'intero sistema (catalogo delle regole, ingestione delle Note, pipeline RAG (\emph{Retrieval-Augmented Generation}), valutazione overnight, snapshot dei risultati canonici, demo Streamlit per la revisione clinica) è riproducibile end-to-end dal repository, con tempi di esecuzione documentati e checkpoint che permettono di riprendere una run interrotta dal punto esatto di interruzione.

\section{Limitazioni del sistema}
\label{sec:concl-limitazioni}

Le limitazioni che il sistema porta con sé al termine del lavoro sono distinte da quelle del framework di valutazione discusse in \Cref{sec:limitazioni}: qui si discute il prodotto, non la misura. La prima e più evidente è la copertura: solo quattro Note AIFA su un catalogo italiano di oltre cinquanta Note vigenti, scelte per eterogeneità clinica e progressione di complessità logica ma non per esaustività di dominio. Il sistema non è quindi un CDSS (\emph{Clinical Decision Support System}) per la rimborsabilità italiana, ma una dimostrazione di fattibilità su un sottoinsieme rappresentativo. La seconda limitazione è il modello generativo: \llmmodel, quantizzato a 4 bit per girare su hardware commodity, è un compromesso fra latenza interattiva e qualità della spiegazione; un modello più grande o un fine-tuning di dominio migliorerebbero presumibilmente la copertura QAEval e l'entailment NLI (\emph{Natural Language Inference}), ma cambierebbero l'ipotesi di deployment locale che è parte costitutiva della tesi. La terza limitazione è l'assenza di integrazione con sistemi gestionali reali: il sistema riceve un profilo paziente già strutturato in formato JSON, non un estratto da cartella clinica elettronica (\emph{Electronic Health Record}, EHR); la trascrizione del profilo dall'EHR alla \texttt{CDSSRequest} è oggi una responsabilità del client e nella pratica clinica richiederebbe un adattatore dedicato per ogni sistema gestionale ospedaliero italiano. La quarta limitazione, infine, è clinica: il sistema non è stato validato in studio prospettico su pazienti reali, e il valore di Macro F1 pari a $1{,}0000$ misura una coerenza interna fra codifica delle regole e annotazione dei casi di riferimento, non una validità clinica esterna. Le implicazioni cliniche del lavoro restano modeste per costruzione; quelle ingegneristiche, come il delta di RQ5 indica, lo sono meno.

\section{Lavori futuri}
\label{sec:concl-futuri}

I naturali sviluppi del lavoro si articolano lungo tre direzioni. La prima è l'estensione della copertura a un sottoinsieme più ampio di Note AIFA: il framework è progettato per accogliere nuove regole con sforzo proporzionale alla complessità della singola Nota, non al numero totale di Note nel sistema. La seconda è il fine-tuning dell'LLM esplicativo su un piccolo insieme di spiegazioni cliniche di qualità: il quadro di RQ3 e RQ6 lascia intuire che le metriche di fedeltà semantica potrebbero migliorare significativamente con un modello adattato al dominio prescrittivo italiano, e l'esperimento sarebbe a basso costo computazionale se condotto in LoRA (\emph{Low-Rank Adaptation}) su Llama 3.1 8B. La terza direzione è la valutazione esterna su un insieme di casi annotato da un panel di prescrittori indipendenti: è l'unico modo per sciogliere la parziale auto-referenza dell'insieme di riferimento ed è la condizione necessaria per qualunque ipotesi di uso clinico del sistema oltre l'ambito sperimentale.

Sul piano della demo, il pannello \emph{what-if} che consente al medico di esplorare interattivamente l'effetto di variazioni dei parametri clinici sulla decisione è implementato e attivo nel flusso principale dell'applicazione; nella versione corrente esso riesegue l'intera pipeline a ogni modifica, e un'ottimizzazione naturale consisterebbe nel ricalcolare la sola spiegazione quando la decisione simbolica resta invariata, riducendo la latenza interattiva. La funzionalità è un ausilio utile soprattutto in fase di formazione del medico, più che di prescrizione operativa.
